\chapter*{KATA PENGANTAR}
\addcontentsline{toc}{chapter}{KATA PENGANTAR}

Puji syukur penulis panjatkan ke hadirat Tuhan Yang Maha Esa atas segala rahmat, hidayah, dan karunia-Nya, sehingga penulis dapat menyelesaikan laporan Tugas Akhir yang berjudul \textbf{"Implementasi Purwarupa Sistem Asesmen Adaptif Kolaboratif Berbasis \textit{Elo Rating} dengan Mitigasi \textit{Overpersonalization}"}. Laporan ini disusun sebagai prasyarat untuk menyelesaikan program pendidikan Sarjana di Program Studi Sistem dan Teknologi Informasi, Sekolah Teknik Elektro dan Informatika, Institut Teknologi Bandung.

Penulis menyadari bahwa penyelesaian laporan ini tidak lepas dari bimbingan, dukungan, dan doa dari berbagai pihak. Oleh karena itu, pada kesempatan ini penulis ingin menyampaikan ucapan terima kasih yang tulus kepada:

\begin{enumerate}
    \item Bapak Ir. Windy Gambetta, MBA., selaku dosen pembimbing yang telah memberikan arahan, bimbingan, masukan yang sangat berharga, serta kesabaran yang luar biasa selama proses pengerjaan Tugas Akhir ini.
    \item Bapak Dr. Riza Satria Perdana, S.T., M.T., selaku dosen wali yang telah membantu dan mendukung penulis dalam mengikuti dan menyelesaikan perkuliahan di Program Studi Sistem dan Teknologi Informasi ITB.
    \item Bapak Ir. I Gusti Bagus Baskara Nugraha, S.T., M.T., Ph.D. dan Ibu Dr. Fetty Fitriyanti Lubis, S.T., M.T., selaku pengampu mata kuliah Tugas Akhir yang telah memberikan informasi terkait proses pengerjaan Tugas Akhir.
    \item Ibu Dr. Fariska Zakhralativa Ruskanda, S.T., M.T., selaku dosen penguji proposal Tugas Akhir yang telah memberikan masukan dan kritik yang membangun untuk perbaikan rencana Tugas Akhir ini.
    \item Kedua orang tua, David Antonius Daliman dan Ika Melania, serta adik tercinta, Latisha Daliman, yang selalu memberikan doa yang tidak pernah putus, dukungan moral, serta motivasi yang menjadi sumber kekuatan bagi penulis sehingga bisa menyelesaikan laporan ini.
    \item Seluruh dosen dan staf Program Studi Sistem dan Teknologi Informasi ITB yang telah membekali penulis dengan ilmu pengetahuan dan pengalaman belajar berarti selama masa perkuliahan yang tentunya menunjang keberhasilan penyusunan Tugas Akhir ini.
    \item Kak Carissa Zahrani Putri, sebagai kakak asuh di perkuliahan yang membantu memberikan arahan dan masukan selama penyusunan Tugas Akhir.
    \item Rekan-rekan mahasiswa angkatan 2022, khususnya teman-teman dari Program Studi Sistem dan Teknologi Informasi ITB dan teman-teman asisten Laboratorium Basis Data, atas diskusi, bantuan, dan kebersamaan yang telah mewarnai perjalanan dan proses penulisan laporan Tugas Akhir ini.
    \item Para partisipan pengujian sistem: Belinda, Darren, Farella, Farhan, Indi, Jelita, Joan, Kenlyn, Nafa, dan Sarah, dari STI ITB angkatan 2023, yang telah meluangkan waktunya untuk membantu proses pengujian dan pemberian masukan pada purwarupa sistem yang dibangun dalam Tugas Akhir ini.
    \item Seluruh pihak yang tidak dapat penulis sebutkan satu per satu, yang telah membantu baik secara langsung maupun tidak langsung dalam penyelesaian laporan ini.
\end{enumerate}

Penulis menyadari bahwa dalam penulisan laporan ini, masih terdapat banyak kekurangan dan keterbatasan. Oleh karena itu, penulis mengharapkan kritik dan saran yang membangun demi perbaikan penulisan karya ilmiah di masa mendatang. Akhir kata, penulis berharap semoga laporan Tugas Akhir ini tidak hanya berhenti di sini sebagai formalitas penyelesaian studi, namun juga dapat memberikan manfaat dan kontribusi positif bagi pengembangan ilmu pengetahuan, khususnya di bidang sistem asesmen adaptif dan teknologi pendidikan.

\begin{flushright}
    Bandung, \today\\
    \vspace{0.8cm}
    Penulis \\
    \vspace{1.2cm}
    Dama Dhananjaya Daliman \\
\end{flushright}
